% !TEX program = xelatex
% AI-effects 프로젝트 --- 직업 규모 분모와 총량 가정이 없는 AI 점유율 지표 A~C와 한국 특화도 지표 D.
% 2026-10-24 재설계: 지표를 AEI 업무용 대화 비중에서 직접 만듦(대화 총량·K 없음, 두 달 = 월별 비중 단순평균), 지표 D 추가.
% 표 본문: results/table/27_summary.tex, 27_top_fs.tex, 27_ksco1.tex, 27_dstats.tex, 27_top_d.tex
%          (code/04_10_analysis_kor_ai_share_index.do가 생성; 지표 계산은 code/03_05_merge_kor_ksco3_ai_usage.do)
% 전체 자료: results/table/27_kor_ai_share_index.xlsx (시트 ksco3, ksco2, ksco1, top_d)
% 관련: paper 19(업무용 대화 비중), paper 25(상대 활용도 RU1·RU2), paper 26(상위 30 소분류)
% 컴파일: xelatex (kotex 필요)
\documentclass[11pt]{article}

\usepackage{fontspec}
\usepackage{kotex}
\usepackage[a4paper,margin=2.2cm]{geometry}
\usepackage{longtable}
\usepackage{booktabs}
\usepackage{array}
\usepackage{amsmath}
\usepackage{xcolor}
\usepackage{hyperref}
\usepackage{enumitem}

\hypersetup{colorlinks=true,linkcolor=blue!50!black,urlcolor=blue!50!black}
\newcolumntype{L}[1]{>{\raggedright\arraybackslash}p{#1}}
\newcolumntype{R}[1]{>{\raggedleft\arraybackslash}p{#1}}
\setlength{\LTleft}{0pt}
\setlength{\LTright}{0pt}
\renewcommand{\arraystretch}{1.15}
\setlength{\tabcolsep}{4pt}

\title{한국 직업별 AI 점유율 지표\\[2mm]
\large --- AEI 업무용 대화 비중으로 만든 지표 A$\sim$C와 한국 특화도 D ---}
\author{AI-effects 프로젝트}
\date{\today}

\begin{document}
\maketitle

\section{왜 필요한가}

paper 25의 상대 활용도 RU1$\cdot$RU2는 직업 규모(취업자 수, 근로시간)로 나눈 강도 지표다. 그래서 취업자가 적은 직업일수록 값이 크게 나올 수 있다(213 생명 및 자연과학 관련 시험원은 취업자 1만 1천 명으로 1위). 취업자 수가 \emph{절대적으로} 얼마나 줄어들 수 있는지를 보려면, 규모로 나누지 않고 한국 전체 중 각 직업의 몫을 재는 지표가 필요하다. paper 26의 상위 30개 선정이 이 방식이었다.

이 지표들은 처음부터 한국 대화 총량과 대화당 메시지 수 $K$가 약분되는 구조였다. 2026년 10월 재설계(총량을 추정하지 않음)에서는 이를 그대로 살려, 총량을 거치지 않고 \textbf{AEI 업무용 대화 비중에서 직접} 만든다. 같은 자료로 한국과 전 세계의 직업 구성을 비교하는 지표 D(한국 특화도)도 더했다.

\section{정의}

소분류 $j$, 월 $t$의 한국 업무용 대화 비중을 $\mathrm{sw}_{jt}$(SOC 업무용 비중을 연계표로 배분한 값, paper 19$\cdot$23), 비중 $\times$ 대화당 절감 시간을 $\mathrm{sv}_{jt}$(SOC 단계에서 곱한 뒤 배분)라 하자. 지표 표본은 한국 비중이 한 달이라도 있는 소분류 129개이고, 월마다 표본 안에서 합이 1이 되게 나눈 뒤 두 달(2026년 4$\cdot$5월)을 단순평균한다. 한 달만 공개된 소분류는 다른 달을 0으로 둔다.
\begin{itemize}[leftmargin=*]
  \item \textbf{A 대화 점유율}: $A_j = 100 \times \tfrac12 \sum_t \mathrm{sw}_{jt} / \sum_k \mathrm{sw}_{kt}$
  \item \textbf{B 절감 시간 점유율}: $B_j = 100 \times \tfrac12 \sum_t \mathrm{sv}_{jt} / \sum_k \mathrm{sv}_{kt}$
  \item \textbf{C 취업자 환산 절감 점유율}: $C_j = 100 \times (B_j/\bar h_j) / \sum_k (B_k/\bar h_k)$. $\bar h_j$는 1인당 월 주업 근로시간(2022년 상반기 지역별고용조사)이다.
  \item \textbf{D 한국 특화도}: $D_j = A_j / A^{\text{GLB}}_j$. $A^{\text{GLB}}_j$는 AEI 전 세계(GLOBAL) 업무용 대화 비중을 같은 식(대분류 업무용 비중, 세부 직업 비중, 같은 연계표, 같은 표본)으로 만든 점유율이다.
  \item (참고) \textbf{A0}: 업무용 보정 없이 모든 use case 대화 비중으로 만든 A. 업무용 보정의 영향을 점검하는 데 쓴다.
\end{itemize}
세 점유율에는 AEI 한국 대화 비중, 대분류 업무용 비중, 연계표, 대화당 절감 시간만 쓰인다. C의 분모 $\bar h_j$는 직업 규모가 아니라 1인당 시간이다. 그래서 같은 절감 시간을 ``몇 명분의 일인가''로 바꿀 뿐 작은 직업을 키우지 않는다.

C는 인원으로 바로 읽을 수 있다. 전국에서 $X$명분의 일이 AI로 대체된다면 소분류 $j$에 배분되는 인원은 $C_j \times X / 100$이다. 직업 간 배분은 $X$와 무관하다. 다만 이 값이 $j$의 취업자 수를 넘으면 안 되므로, 배분 인원이 취업자 수와 같아지는 $X$를 상한으로 점검했다. $X$ 자체(노동 절감의 총량)는 이번 재설계에서 추계하지 않는다.

D는 ``한국 업무용 AI 사용 중 이 직업의 몫이 전 세계에서의 몫의 몇 배인가''다. 1이면 전 세계와 같은 몫, 1보다 크면 한국에서 상대적으로 더 많이 쓰는 직업이다. 직업 구성의 차이만 반영하도록 GLOBAL도 한국과 같은 방식(세부 직업 업무용 비중은 대분류 값을 씀)으로 만들었다. 정의상 GLOBAL 점유율로 가중한 D의 평균은 1이다.

\section{결과}

\subsection{지표 A$\sim$C}

\begin{longtable}{L{5.5cm}R{2.3cm}R{2.3cm}R{2.6cm}R{2.6cm}}
\caption{지표 A$\sim$C의 집중도와 RU1$\cdot$RU2와의 순위상관(소분류 129개)}\label{tab:summary}\\
\toprule
지표 & 상위 10개 합(\%) & 상위 30개 합(\%) & RU1과 순위상관 & RU2와 순위상관 \\
\midrule
\endfirsthead
% 04_10_analysis_kor_ai_share_index.do가 생성. 직접 고치지 말 것.
A 대화 점유율 & 54.7 & 81.3 & 0.81 & 0.78 \\
B 절감 시간 점유율 & 55.5 & 84.1 & 0.84 & 0.86 \\
C 취업자 환산 점유율 & 55.5 & 84.4 & 0.84 & 0.86 \\
(참고) A0 모든 use case 대화 점유율 & 50.7 & 79.8 & 0.74 & 0.70 \\
(참고) 취업자 비중 & 38.3 & 68.2 & -- & -- \\

\bottomrule
\end{longtable}

\begin{longtable}{R{0.6cm}L{4.0cm}R{1.2cm}R{1.3cm}R{1.6cm}R{1.6cm}R{1.3cm}R{1.2cm}R{1.1cm}}
\caption{지표 C 상위 20개 소분류}\label{tab:top}\\
\toprule
순위 & 소분류 & C 점유율(\%) & 1만 명당 배분(명) & A 점유율(순위) & B 점유율(순위) & 취업자(천 명) & RU2 순위 & D \\
\midrule
\endfirsthead
\toprule
순위 & 소분류 & C 점유율(\%) & 1만 명당 배분(명) & A 점유율(순위) & B 점유율(순위) & 취업자(천 명) & RU2 순위 & D \\
\midrule
\endhead
% 04_10_analysis_kor_ai_share_index.do가 생성. 직접 고치지 말 것.
1 & 222 컴퓨터 시스템 및 소프트웨어 전문가 & 11.56 & 1,156 & 9.95 (2) & 12.21 (1) & 399 & 22 & 0.86 \\
2 & 291 작가 및 언론 관련 전문가 & 11.08 & 1,108 & 11.92 (1) & 10.47 (2) & 96 & 5 & 1.30 \\
3 & 372 비서 및 사무 지원 종사자 & 5.81 & 581 & 4.69 (4) & 5.24 (4) & 351 & 31 & 1.12 \\
4 & 225 정보 시스템 및 웹 운영자 & 5.13 & 513 & 8.32 (3) & 5.34 (3) & 79 & 13 & 0.90 \\
5 & 283 상품 기획‧홍보 및 조사 전문가 & 5.06 & 506 & 4.20 (5) & 5.15 (5) & 62 & 9 & 0.99 \\
6 & 332 통계‧데이터 관련 사무원 & 4.20 & 420 & 2.61 (9) & 4.03 (6) & 16 & 2 & 1.16 \\
7 & 213 생명 및 자연과학 관련 시험원 & 3.51 & 351 & 2.79 (8) & 3.62 (8) & 11 & 1 & 1.16 \\
8 & 223 네트워크 및 정보 보안 전문가 & 3.51 & 351 & 3.84 (6) & 3.74 (7) & 33 & 6 & 0.89 \\
9 & 293 연극‧영화 및 영상 전문가 & 2.91 & 291 & 2.43 (11) & 2.85 (10) & 120 & 27 & 1.27 \\
10 & 211 생명 및 자연과학 관련 전문가 & 2.78 & 278 & 2.04 (12) & 2.89 (9) & 76 & 17 & 1.07 \\
11 & 264 문리‧기술 및 예능 강사 & 2.74 & 274 & 1.85 (14) & 2.07 (13) & 585 & 50 & 1.12 \\
12 & 292 학예사‧사서 및 기록물 관리사 & 2.37 & 237 & 3.81 (7) & 2.29 (12) & 31 & 10 & 1.32 \\
13 & 269 기타 교육 전문가 & 2.33 & 233 & 1.41 (16) & 2.02 (14) & 46 & 15 & 1.24 \\
14 & 282 금융 및 보험 전문가 & 2.30 & 230 & 1.98 (13) & 2.48 (11) & 47 & 16 & 1.12 \\
15 & 340 금융 사무 종사자 & 1.73 & 173 & 2.61 (10) & 1.78 (16) & 287 & 48 & 0.96 \\
16 & 224 데이터 전문가 & 1.72 & 172 & 1.10 (21) & 1.82 (15) & 25 & 11 & 0.80 \\
17 & 235 기계‧로봇공학 기술자 및 시험원 & 1.59 & 159 & 1.71 (15) & 1.70 (17) & 122 & 35 & 1.16 \\
18 & 122 마케팅 및 광고‧홍보 관리자 & 1.44 & 144 & 1.01 (25) & 1.50 (18) & 9 & 4 & 0.75 \\
19 & 261 대학교수 및 강사 & 1.33 & 133 & 1.01 (26) & 1.12 (24) & 113 & 36 & 1.18 \\
20 & 131 연구‧교육 및 법률 관련 관리자 & 1.28 & 128 & 1.39 (17) & 1.35 (19) & 50 & 24 & 0.87 \\

\bottomrule
\end{longtable}

\begin{longtable}{L{5.2cm}R{1.6cm}R{1.8cm}R{1.8cm}R{1.8cm}R{1.8cm}R{1.4cm}}
\caption{대분류별 점유율(\%, 소분류 합)과 D}\label{tab:ksco1}\\
\toprule
대분류 & 소분류 수 & A 대화 & B 절감 시간 & C 취업자 환산 & 취업자 비중 & D \\
\midrule
\endfirsthead
% 04_10_analysis_kor_ai_share_index.do가 생성. 직접 고치지 말 것.
1 관리자 & 16 & 6.9 & 8.9 & 8.4 & 1.7 & 0.83 \\
2 전문가 및 관련 종사자 & 45 & 70.5 & 72.2 & 71.9 & 23.6 & 1.02 \\
3 사무 종사자 & 13 & 14.6 & 14.4 & 15.0 & 19.8 & 1.06 \\
4 서비스 종사자 & 13 & 1.1 & 0.5 & 0.6 & 11.4 & 1.02 \\
5 판매 종사자 & 5 & 3.2 & 1.3 & 1.2 & 10.8 & 0.79 \\
6 농림∙어업 숙련 종사자 & 4 & 0.0 & 0.0 & 0.0 & 5.9 & 1.67 \\
7 기능원 및 관련 기능 종사자 & 14 & 1.9 & 1.4 & 1.5 & 7.0 & 0.97 \\
8 장치∙기계 조작 및 조립 종사자 & 10 & 0.9 & 1.0 & 0.9 & 6.6 & 0.88 \\
9 단순노무 종사자 & 9 & 0.9 & 0.4 & 0.5 & 13.1 & 0.97 \\

\bottomrule
\end{longtable}

\begin{itemize}[leftmargin=*]
  \item \textbf{재설계 전과 거의 같다.} 총량을 거치던 예전 계산과 비교하면, 두 달 모두 공개된 소분류 122개의 A$\cdot$B는 상대 차이 0.04\% 이내다. 한 달만 공개된 하위권 7개만 다른 달을 0으로 두면서 값이 절반이 되었다. 순위상관은 0.9999(A), 0.9998(B)이고, 상위 10$\cdot$30개 합도 A 상위 30개(81.2\% $\to$ 81.3\%) 외에는 소수 첫째 자리까지 같다.
  \item \textbf{B와 C는 거의 같다}(순위상관 0.999). 1인당 근로시간 차이가 작아 보정 효과가 작기 때문이다. A와 C의 순위상관은 0.96이다. 대화당 절감 시간이 긴 직업(332 통계$\cdot$데이터 사무원, 264 문리$\cdot$기술 및 예능 강사 등)은 C에서 순위가 오르고, 짧은 직업(225 정보 시스템 및 웹 운영자, 292 학예사$\cdot$사서)은 내려간다.
  \item \textbf{강도 지표와는 다른 직업을 가리킨다.} C 상위 30개 중 RU2 상위 30개와 겹치는 것은 19개다. 264 문리$\cdot$기술 및 예능 강사는 C 11위지만 RU2는 50위다(취업자 58만 5천 명). 반대로 213 생명 및 자연과학 관련 시험원은 RU2 1위지만 C로는 7위다.
  \item \textbf{집중도가 높다.} 상위 10개 소분류가 C의 55.5\%, 상위 30개가 84.4\%를 차지한다. 취업자가 가장 많은 소분류 10개$\cdot$30개의 취업자 비중 합(38.3\%, 68.2\%)보다 높다.
  \item \textbf{전문가 대분류에 몰린다.} 전문가 및 관련 종사자는 취업자의 23.6\%지만 C의 71.9\%를 차지한다. 사무 종사자는 15.0\%(취업자 19.8\%)다. 서비스$\cdot$판매$\cdot$단순노무는 합쳐 2.3\%다.
  \item \textbf{상한}: 전국 약 32만 7천 명분까지는 어느 소분류도 배분 인원이 자기 취업자 수를 넘지 않는다. 가장 먼저 상한에 닿는 것은 213 생명 및 자연과학 관련 시험원이다.
  \item \textbf{업무용 보정(A 대 A0)}: 업무용 보정을 빼면 집중도가 조금 낮아진다(상위 30개 합 81.3\% $\to$ 79.8\%). A와 A0의 순위상관은 0.957이고 상위 30개 중 27개가 겹친다. 업무용 비중이 낮은 교육$\cdot$판매$\cdot$보건 쪽 직업이 A0에서 커진다(paper 19).
\end{itemize}

\subsection{지표 D: 한국 특화도}

\begin{longtable}{L{9.5cm}R{3.5cm}}
\caption{지표 D의 분포(소분류)}\label{tab:dstats}\\
\toprule
항목 & 값 \\
\midrule
\endfirsthead
% 04_10_analysis_kor_ai_share_index.do가 생성. 직접 고치지 말 것.
D 중앙값(129개) & 0.99 \\
D 10분위 / 90분위(129개) & 0.68 / 1.28 \\
D 최솟값 / 최댓값(129개) & 0.02 / 2.41 \\
D 10분위 / 90분위(A $\ge$ 0.1\% 소분류 83개) & 0.75 / 1.21 \\
D $>$ 1인 소분류 수(한국 A 합, \%) & 55 (49.7) \\
한국 A로 가중한 D 평균 & 1.03 \\
D와 A의 Spearman 순위상관 & 0.10 \\
D와 RU1의 Spearman 순위상관 & -0.02 \\
A와 A0의 Spearman 순위상관 & 0.957 \\
A와 A0의 상위 30 겹침(개) & 27 \\

\bottomrule
\end{longtable}

\begin{longtable}{L{6.4cm}R{1.4cm}R{2.4cm}R{2.4cm}R{2.0cm}}
\caption{D 상위$\cdot$하위 15개 소분류(A $\ge$ 0.1\%인 83개 중)}\label{tab:topd}\\
\toprule
소분류 & D & A 점유율(\%, 순위) & GLOBAL 점유율(\%) & 취업자(천 명) \\
\midrule
\endfirsthead
\toprule
소분류 & D & A 점유율(\%, 순위) & GLOBAL 점유율(\%) & 취업자(천 명) \\
\midrule
\endhead
\multicolumn{5}{l}{\textit{한국이 더 많이 쓰는 직업(D 상위 15)}} \\
% 04_10_analysis_kor_ai_share_index.do가 생성. 직접 고치지 말 것.
782 건설 관련 기능 종사자 & 1.52 & 0.13 (78) & 0.08 & 295 \\
292 학예사‧사서 및 기록물 관리사 & 1.32 & 3.81 (7) & 2.89 & 31 \\
291 작가 및 언론 관련 전문가 & 1.30 & 11.92 (1) & 9.17 & 96 \\
293 연극‧영화 및 영상 전문가 & 1.27 & 2.43 (11) & 1.92 & 120 \\
312 인사 및 교육·훈련 사무원 & 1.25 & 0.45 (43) & 0.36 & 203 \\
269 기타 교육 전문가 & 1.24 & 1.41 (16) & 1.14 & 46 \\
133 보건 및 사회복지 관련 관리자 & 1.22 & 0.26 (57) & 0.21 & 11 \\
233 금속‧재료공학 기술자 및 시험원 & 1.22 & 0.22 (62) & 0.18 & 12 \\
262 학교 교사 & 1.21 & 0.14 (74) & 0.12 & 422 \\
236 소방‧방재‧산업안전 기술자 및 시험원 & 1.19 & 0.69 (33) & 0.58 & 120 \\
212 인문 및 사회과학 전문가 & 1.19 & 1.26 (19) & 1.07 & 22 \\
239 기타 공학 전문가 및 관련 종사자 & 1.18 & 1.07 (23) & 0.90 & 110 \\
237 환경공학‧가스‧에너지 기술자 및 시험원 & 1.18 & 0.37 (48) & 0.31 & 40 \\
261 대학교수 및 강사 & 1.18 & 1.01 (26) & 0.86 & 113 \\
254 종교 관련 종사자 & 1.17 & 0.19 (67) & 0.16 & 116 \\
\midrule
224 데이터 전문가 & 0.80 & 1.10 (21) & 1.37 & 25 \\
352 정부 행정 사무 종사자 & 0.79 & 0.63 (36) & 0.79 & 523 \\
132 보험 및 금융 관리자 & 0.77 & 0.22 (61) & 0.28 & 46 \\
510 영업 종사자 & 0.76 & 0.92 (28) & 1.22 & 609 \\
532 방문 및 노점 판매 종사자 & 0.75 & 0.12 (81) & 0.16 & 85 \\
122 마케팅 및 광고‧홍보 관리자 & 0.75 & 1.01 (25) & 1.35 & 9 \\
295 디자이너 & 0.75 & 1.19 (20) & 1.59 & 266 \\
531 통신 관련 판매 종사자 & 0.73 & 0.22 (60) & 0.30 & 255 \\
272 정부 행정 전문가 및 관련 종사자 & 0.73 & 0.36 (49) & 0.49 & 31 \\
139 기타 전문 서비스 관리자 & 0.72 & 0.24 (59) & 0.34 & 21 \\
953 판매 관련 단순 종사자 & 0.72 & 0.13 (79) & 0.17 & 178 \\
522 상품 대여 종사자 & 0.70 & 1.07 (22) & 1.55 & 30 \\
151 판매 및 운송 관리자 & 0.69 & 0.34 (51) & 0.49 & 66 \\
281 인사 및 경영 전문가 & 0.68 & 0.52 (42) & 0.76 & 86 \\
152 고객 서비스 관리자 & 0.64 & 0.20 (64) & 0.32 & 15 \\

\bottomrule
\multicolumn{5}{l}{\footnotesize 가운데 구분선 아래는 D 하위 15(한국이 덜 쓰는 직업).}
\end{longtable}

\begin{itemize}[leftmargin=*]
  \item \textbf{한국의 직업 구성은 전 세계와 대체로 비슷하다.} D의 중앙값은 0.99이고, 점유율이 0.1\% 이상인 83개 소분류에서 10$\sim$90분위는 0.75$\sim$1.21이다. 극단값(0.02, 2.41)은 점유율이 0.01\% 안팎인 소분류(451 조리사, 864 전기$\cdot$전자 부품 조립원, 611$\cdot$620$\cdot$630 농림어업)에서 나온다. 이 때문에 목록은 A $\ge$ 0.1\%로 제한했다.
  \item \textbf{한국에서 상대적으로 더 많이 쓰는 직업}: 1위는 782 건설 관련 기능 종사자(1.52)지만 점유율이 0.13\%로 작다. 그 밖에는 글쓰기$\cdot$기록(292 학예사$\cdot$사서 1.32, 291 작가$\cdot$언론 1.30, 293 연극$\cdot$영화$\cdot$영상 1.27), 교육(269 기타 교육 전문가, 262 학교 교사, 261 대학교수), 공학 기술자(233, 236, 237, 239)다. 대분류 단위에서는 사무(1.06)와 전문가(1.02)가 1보다 조금 크다.
  \item \textbf{한국에서 상대적으로 덜 쓰는 직업}: 데이터$\cdot$마케팅(224 데이터 전문가 0.80, 122 마케팅$\cdot$광고 관리자 0.75, 295 디자이너 0.75), 영업$\cdot$판매(510 영업 0.76, 531 통신 판매 0.73, 522 상품 대여 0.70), 관리$\cdot$경영(132, 139, 151, 152, 281)이다. 대분류로는 판매(0.79)와 관리자(0.83)가 낮다.
  \item \textbf{D는 사용 강도와 다른 축이다.} D와 A의 순위상관은 0.10, D와 RU1은 $-0.02$다. 한국에서 많이 쓰는 직업이 곧 한국이 특히 많이 쓰는 직업은 아니다. 대분류 단위 차이(paper 19 표 3: 교육$\cdot$도서관 1.29, 예술$\cdot$미디어 1.19 대 경영 0.81, 판매 0.77)가 소분류 D에 그대로 옮겨진 부분이 크다.
\end{itemize}

\section{해석상 주의}

\begin{enumerate}[leftmargin=*]
  \item \textbf{총량은 정하지 않는다}: 점유율은 합이 1이라 배분만 알려 준다. 전국 감소 규모 $X$와 시점 간 증가는 따로 정해야 하며, 이번 재설계에서는 추계하지 않는다.
  \item \textbf{상한형 배분}: C는 절감 시간이 모두 인원 감소로 이어진다고 볼 때의 배분이다. 절감 시간은 산출 증가로 흡수될 수도 있다. 자동화(automation) 대화로 좁히면 더 보수적인 지표가 된다(AEI GLOBAL task의 협업 유형).
  \item \textbf{기존 가정}: AEI 직업은 대화 과업의 직업이다. 한국 업무용 AI 사용의 직업 구성이 Claude와 같다고 보고(paper 19의 B1), 연계표 배분(paper 21)을 따른다. 15-1299나 Data Scientists의 배분이 점유율 순위에 직접 영향을 준다.
  \item \textbf{D의 해석}: D는 Claude.ai 안에서 한국과 전 세계의 직업 구성을 비교한다. 한국 Claude 이용자의 구성(직업$\cdot$연령$\cdot$언어)이 전 세계와 다르면 그 차이도 D에 들어간다. 세부 직업 업무용 비중은 두 지역 모두 대분류 값을 썼다(가정 B3; GLOBAL에서 B3을 점검한 결과는 paper 19).
  \item \textbf{$\bar h_j$의 시점}: paper 25와 같이 2022년 상반기 값이다. C는 B와 거의 같아 시점 선택의 영향은 작다.
\end{enumerate}

\section{산출 파일}

\begin{itemize}[leftmargin=*]
  \item \texttt{code/03\_05\_merge\_kor\_ksco3\_ai\_usage.do}: 지표 A$\sim$D와 A0를 계산해 \texttt{data/proc/merged/kor\_ai\_usage\_ksco\{3,2,1\}}에 저장한다.
  \item \texttt{code/04\_10\_analysis\_kor\_ai\_share\_index.do}: 이 노트의 표를 만든다.
  \item \texttt{results/table/27\_kor\_ai\_share\_index.xlsx}: 시트 \texttt{ksco3}(소분류 129개의 A$\sim$D, A0, GLOBAL 점유율, 순위, 1만 명당 배분, 상한, RU1$\cdot$RU2), \texttt{ksco2}(중분류 합, 50개), \texttt{ksco1}(대분류 합), \texttt{top\_d}(D 상위$\cdot$하위 15).
  \item \texttt{results/table/27\_summary.tex}, \texttt{27\_top\_fs.tex}, \texttt{27\_ksco1.tex}, \texttt{27\_dstats.tex}, \texttt{27\_top\_d.tex}: 표~\ref{tab:summary}--\ref{tab:topd}.
\end{itemize}

\end{document}
